\documentclass[aspectratio=169,10pt]{beamer}

\usetheme{Madrid}
\usecolortheme{default}

\usepackage{fontspec}
\usepackage{amsmath}
\usepackage{booktabs}
\usepackage{tabularx}
\usepackage{array}
\usepackage{ragged2e}
\usepackage{makecell}

\setsansfont{Arial}
\setmonofont{Consolas}

\definecolor{MadridBlue}{RGB}{36,74,122}
\definecolor{SoftBlue}{RGB}{232,240,248}
\definecolor{SoftGray}{RGB}{244,245,247}
\definecolor{DarkGray}{RGB}{70,74,82}
\definecolor{GoodGreen}{RGB}{24,112,72}
\definecolor{WarnRed}{RGB}{165,52,52}

\setbeamercolor{structure}{fg=MadridBlue}
\setbeamercolor{alerted text}{fg=WarnRed}
\setbeamertemplate{navigation symbols}{}
\setbeamersize{text margin left=0.55cm,text margin right=0.55cm}

\newcommand{\key}[1]{\textcolor{MadridBlue}{\textbf{#1}}}
\newcommand{\risk}[1]{\textcolor{WarnRed}{\textbf{#1}}}
\newcommand{\good}[1]{\textcolor{GoodGreen}{\textbf{#1}}}
\newcommand{\source}[1]{\vfill{\scriptsize\textcolor{DarkGray}{#1}}}
\newcommand{\tightitems}{\setlength\itemsep{0.11cm}\setlength\parskip{0pt}}

\title[Dissertation Essays]{Dissertation Essays: Supervisor Meeting}
\author{}
\institute{}
\date{8 September 2026}

\begin{document}

% -----------------------------------------------------------------------------
\begin{frame}[t]{Essay 1: current paper in one slide}
\begin{columns}[T,onlytextwidth]
  \begin{column}{0.54\textwidth}
    \begin{block}{Question and setting}
      Do firm website restrictions on search and discovery crawlers change market liquidity around mandatory disclosures?
    \end{block}
    \footnotesize
    \begin{itemize}\tightitems
      \item Treatment: the first restriction of any of 31 search/discovery crawlers in archived monthly \texttt{robots.txt} files.
      \item Not an AI paper: the treatment covers the full search/discovery layer, including conventional search bots.
      \item Design: staggered adoption from 2021 to 2025; 27,327 10-K/10-Q filing events for 1,908 firms.
      \item Main outcome: filing-window log Amihud, measured over $[-1,+1]$, relative to the firm's local benchmark $[-60,-10]$.
    \end{itemize}
  \end{column}
  \begin{column}{0.43\textwidth}
    \begin{alertblock}{Current headline result}
      BJS estimate: \textbf{+0.0134}, $p=0.010$. The TWFE estimate is \textbf{+0.0109}.
    \end{alertblock}
    \footnotesize
    \begin{itemize}\tightitems
      \item The filing window does not become less liquid in levels; its local benchmark improves faster.
      \item Quoted spreads are null, dollar volume rises, and return variance rises.
      \item Effects are concentrated among smaller, low-institutional-ownership, and low-coverage firms.
    \end{itemize}
    \begin{exampleblock}{Current interpretation}
      \footnotesize
      The evidence points toward price impact/depth or trader composition, not wider quoted spreads or trading withdrawal.
    \end{exampleblock}
  \end{column}
\end{columns}
\end{frame}

% -----------------------------------------------------------------------------
\begin{frame}[t]{Feedback from the committee and our discussion}
\vspace{0.20cm}
\begin{columns}[T,onlytextwidth]
  \begin{column}{0.31\textwidth}
    \begin{block}{1. Measurement paper}
      \small
      Candidate contribution: the first systematic evidence that firms \textbf{actively make their own public information harder for machines to discover}.

      \vspace{0.12cm}
      This requires a first stage in actual discoverability, not only a stated \texttt{robots.txt} policy.
    \end{block}
  \end{column}
  \begin{column}{0.31\textwidth}
    \begin{block}{2. Mechanism}
      \small
      \textbf{Amihud is not information deterioration.}

      \vspace{0.12cm}
      Separate lower information acquisition, more informed trading, noise-trader exit, and changes in disclosure content.
    \end{block}
  \end{column}
  \begin{column}{0.31\textwidth}
    \begin{block}{3. More output variables}
      \small
      Expand the outcome set, but require each variable to adjudicate among mechanisms.

      \vspace{0.12cm}
      Priority families: price response, information timing, investor breadth, monitoring, and disclosure behavior.
    \end{block}
  \end{column}
\end{columns}
\end{frame}

% -----------------------------------------------------------------------------
\begin{frame}[t]{The dissertation studies effective, not merely legal, publicity}
\begin{block}{Common theoretical primitive}
\centering
\Large
$\text{Effective public information}
=
\text{legal availability}
\times
\text{discoverability}
\times
\text{processing capacity}$
\end{block}

\vspace{0.18cm}
\small
Information can be legally public and still behave like private information if outsiders cannot find it cheaply and repeatedly. The three essays study a different equilibrium margin of this same friction:

\vspace{0.18cm}
\begin{columns}[T,onlytextwidth]
  \begin{column}{0.31\textwidth}
    \begin{exampleblock}{Trading margin}
      Who trades, how much depth they supply, and how strongly prices move per unit of volume.
    \end{exampleblock}
  \end{column}
  \begin{column}{0.31\textwidth}
    \begin{exampleblock}{Monitoring margin}
      Who continuously reads firm information and how managers respond when monitoring becomes costly.
    \end{exampleblock}
  \end{column}
  \begin{column}{0.31\textwidth}
    \begin{exampleblock}{Intermediation margin}
      Whether open-web discovery substitutes for or complements traditional human intermediaries.
    \end{exampleblock}
  \end{column}
\end{columns}

\source{Conceptual anchors: Grossman and Stiglitz (1980), Merton (1987), Kyle (1985), and Dugast and Foucault (2018).}
\end{frame}

% -----------------------------------------------------------------------------
\begin{frame}[t]{Three essays share a friction, not an outcome swap}
\scriptsize
\renewcommand{\arraystretch}{1.20}
\setlength{\tabcolsep}{0.11cm}
\begin{tabularx}{\textwidth}{>{\bfseries\RaggedRight\arraybackslash}p{2.25cm} >{\RaggedRight\arraybackslash}X >{\RaggedRight\arraybackslash}X >{\RaggedRight\arraybackslash}X}
\toprule
 & \textbf{Essay 1} & \textbf{Essay 2} & \textbf{Essay 3} \\
\midrule
Working title & Crawler-Access Restrictions and Disclosure-Window Liquidity & Turning Off the Machines & Searchability and Information Intermediary Substitution \\
Literature home & Market microstructure & Corporate governance & Information intermediaries \\
Economic margin & Trading and price impact & External monitoring & Substitution vs. complementarity \\
Main shock & Firm website crawler restriction & Firm website crawler restriction & Broker closure or research-department shutdown \\
Role of discoverability & Treatment and measured friction & Treatment plus monitoring exposure & Predetermined moderator, not treatment \\
Primary object & Liquidity and the timing of price learning & Opportunistic managerial behavior & Damage caused by exogenous analyst loss \\
\bottomrule
\end{tabularx}

\vspace{0.20cm}
\begin{alertblock}{Important correction}
\footnotesize
Essay 2 still uses the robots restriction shock. It is independent only if it establishes a distinct monitoring channel and a prespecified governance outcome. Essay 3 genuinely changes the treatment.
\end{alertblock}
\end{frame}

% -----------------------------------------------------------------------------
\begin{frame}[t]{Essay 1 measures a firm choice over discoverability}
\begin{columns}[T,onlytextwidth]
  \begin{column}{0.48\textwidth}
    \begin{block}{What the measure captures}
    \small
    \begin{itemize}\tightitems
      \item A monthly policy governing whether named crawlers may request content from the firm's website.
      \item A firm's ability to change the machine accessibility of information that remains visible to a human visitor.
      \item A broad search/discovery infrastructure choice, not a narrow response to generative AI.
    \end{itemize}
    \end{block}
  \end{column}
  \begin{column}{0.48\textwidth}
    \begin{block}{What it does not yet prove}
    \small
    \begin{itemize}\tightitems
      \item That major crawlers obeyed the restriction at the relevant time.
      \item That indexed pages, ranking, or discovery latency actually changed.
      \item That managers personally intended the change rather than inheriting a vendor or security template.
    \end{itemize}
    \end{block}
  \end{column}
\end{columns}

\begin{exampleblock}{Careful contribution language}
Until intent and the actual discovery first stage are validated, the defensible statement is: \textbf{firms' websites impose restrictions that reduce permitted machine access}. ``Firms deliberately hide'' is the stronger hypothesis to be tested.
\end{exampleblock}
\end{frame}

% -----------------------------------------------------------------------------
\begin{frame}[t]{Essay 1 uses filing events to isolate a high-information window}
\begin{columns}[T,onlytextwidth]
  \begin{column}{0.52\textwidth}
    \begin{block}{Empirical design}
    \small
    \begin{itemize}\tightitems
      \item Unit: firm filing event for 10-K and 10-Q disclosures.
      \item Treatment date: first observed restriction of any focal search/discovery crawler.
      \item Comparison: matched never-restricting firms and not-yet-treated observations.
      \item Estimators: BJS imputation as the main design; conventional TWFE as a transparent comparison.
      \item Outcome: abnormal filing-window log Amihud relative to a firm-local pre-filing benchmark.
    \end{itemize}
    \end{block}
  \end{column}
  \begin{column}{0.44\textwidth}
    \begin{block}{Why use a local benchmark?}
    \small
    The estimand compares the filing window with the same firm's recent normal trading environment. This removes much of the persistent cross-sectional difference in baseline liquidity.
    \end{block}
    \begin{alertblock}{Interpretation constraint}
    \small
    A positive abnormal measure may arise because the event window worsens, because normal days improve, or both. The two components must be shown separately.
    \end{alertblock}
  \end{column}
\end{columns}
\end{frame}

% -----------------------------------------------------------------------------
\begin{frame}[t]{The main effect is relative deterioration, not a level collapse}
\begin{columns}[T,onlytextwidth]
  \begin{column}{0.37\textwidth}
    \begin{alertblock}{Main estimate}
      \centering
      \LARGE \textbf{+0.0134}

      \normalsize BJS abnormal log Amihud

      \small $p=0.010$
    \end{alertblock}
    \begin{block}{Concentration}
      \footnotesize The effect is stronger for small firms, firms with low institutional ownership, and firms with little analyst coverage.
    \end{block}
  \end{column}
  \begin{column}{0.59\textwidth}
    \begin{block}{What moves underneath the ratio?}
    \footnotesize
    \begin{itemize}\tightitems
      \item Filing-window Amihud does not rise materially in levels.
      \item The non-event benchmark improves faster after restrictions.
      \item Dollar volume rises rather than falls.
      \item Quoted spreads do not widen.
      \item Return variance rises around filings.
      \item Post-filing drift and delayed-response share are currently null.
    \end{itemize}
    \end{block}
    \begin{exampleblock}{Most accurate current reading}
      \footnotesize The filing window becomes relatively more price-impact intensive. The evidence does not yet distinguish lower depth, concentrated information arrival, and trader composition.
    \end{exampleblock}
  \end{column}
\end{columns}
\end{frame}

% -----------------------------------------------------------------------------
\begin{frame}[t,shrink=3]{Amihud can increase for economically different reasons}
\begin{block}{The measure}
\centering
\Large $\text{Amihud}_{it}=10^{6}\,\dfrac{|r_{it}|}{\text{DollarVolume}_{it}}$
\end{block}

\scriptsize
\renewcommand{\arraystretch}{1.14}
\setlength{\tabcolsep}{0.10cm}
\begin{tabularx}{\textwidth}{>{\bfseries\RaggedRight\arraybackslash}p{2.50cm} >{\RaggedRight\arraybackslash}X >{\RaggedRight\arraybackslash}X}
\toprule
Channel & Why Amihud rises & Additional implication \\
\midrule
Lower market depth & A given order imbalance moves price more & Effective-spread or price-impact proxies rise; breadth may fall \\
More informative trades & Returns respond more strongly to fundamentals & Return-news covariance rises, even if liquidity is unchanged \\
Noise traders exit & Denominator falls because uninformed volume disappears & Volume and participation breadth fall \\
Information becomes concentrated & More of the quarterly news arrives at the filing & Pre-filing learning falls; filing response rises \\
Local benchmark changes & Normal-day liquidity improves faster than event liquidity & Abnormal measure rises without an event-window level increase \\
\bottomrule
\end{tabularx}

\begin{alertblock}{Therefore}
\footnotesize The paper cannot infer ``information worsens'' from Amihud. It must identify the component and price-learning pattern that changed.
\end{alertblock}
\end{frame}

% -----------------------------------------------------------------------------
\begin{frame}[t]{Essay 1 needs a mechanism horse race, not a larger outcome menu}
\scriptsize
\renewcommand{\arraystretch}{1.25}
\setlength{\tabcolsep}{0.08cm}
\begin{tabularx}{\textwidth}{>{\bfseries\RaggedRight\arraybackslash}p{2.35cm} >{\RaggedRight\arraybackslash}X >{\RaggedRight\arraybackslash}X >{\RaggedRight\arraybackslash}X >{\RaggedRight\arraybackslash}X}
\toprule
Mechanism & Volume / breadth & Spread / depth & Fundamental response & Post-filing learning \\
\midrule
Information loss & May fall & May worsen & Full-cycle response falls & Drift or delayed response rises \\
Temporal concentration & Need not fall & Event price impact may rise & Pre-filing down; filing up; total unchanged & Little residual learning \\
Informed-trader composition & Volume may rise or stay flat & Price impact rises more than quoted spread & Response per trade rises & No necessary drift \\
Noise-trader exit & Volume and breadth fall & Price impact may rise & Fundamental mapping need not change & No sharp prediction \\
Disclosure-content change & Depends on news quantity and complexity & Depends on content & Changes are explained by filing news & Depends on omitted content \\
\bottomrule
\end{tabularx}

\vspace{0.20cm}
\begin{exampleblock}{How the current facts narrow the race}
Higher dollar volume and null quoted spreads make simple trading withdrawal less plausible. Higher return variance with no post-filing drift is consistent with either concentrated information arrival or a change in price impact. The next tests must distinguish these two.
\end{exampleblock}

\source{Theoretical motivation: Kyle (1985); Merton (1987); Davila and Parlatore (2026) on the non-monotone relation between trading costs and price informativeness.}
\end{frame}

% -----------------------------------------------------------------------------
\begin{frame}[t]{``Where Prices Learn'' belongs inside Essay 1}
\begin{block}{Its role in the paper}
\large
It explains \textbf{why the same mandatory filing may generate greater price movement per unit of trading after search access is restricted}.
\end{block}

\small
\begin{columns}[T,onlytextwidth]
  \begin{column}{0.47\textwidth}
    \begin{exampleblock}{Before restriction}
      Search and discovery tools expose continuous website signals: IR updates, product news, hiring, operational announcements, and changes to firm pages.

      Some quarterly information may therefore enter prices before the filing date.
    \end{exampleblock}
  \end{column}
  \begin{column}{0.47\textwidth}
    \begin{alertblock}{After restriction}
      The cost of finding continuous signals rises. The mandatory SEC filing remains accessible and may become a larger information reset.

      Price learning may shift from dispersed pre-filing days to the filing window.
    \end{alertblock}
  \end{column}
\end{columns}

\vspace{0.12cm}
\textbf{This is a mechanism for the Essay 1 market result.} With the same restriction events, firms, filing cycle, and return data, it would look like salami slicing if presented as a separate essay.
\end{frame}

% -----------------------------------------------------------------------------
\begin{frame}[t]{The relevant question is when and through which channel prices learn}
\begin{center}
\begin{tabular}{ccc}
\colorbox{SoftBlue}{\parbox{0.27\textwidth}{\centering\textbf{PRE-FILING}\\Continuous web signals\\Anticipatory incorporation}} &
$\Longrightarrow$ &
\colorbox{SoftBlue}{\parbox{0.27\textwidth}{\centering\textbf{FILING}\\Mandatory disclosure\\Scheduled information reset}} \\
&&\\[-0.18cm]
&& $\Downarrow$ \\
&&\\[-0.18cm]
&& \colorbox{SoftGray}{\parbox{0.27\textwidth}{\centering\textbf{POST-FILING}\\Residual learning\\Delay or correction?}}
\end{tabular}
\end{center}

\begin{block}{Terminology}
\small
``Where'' does not mean a physical trading venue. It means \textbf{where in the information cycle} and \textbf{through which disclosure channel} fundamental news enters prices. A more precise section title is: \emph{When and Through Which Channel Do Prices Learn?}
\end{block}

\end{frame}

% -----------------------------------------------------------------------------
\begin{frame}[t,shrink=2]{Price learning requires a fundamental anchor}
\begin{columns}[T,onlytextwidth]
  \begin{column}{0.48\textwidth}
    \begin{block}{Primary measures}
    \small
    \begin{itemize}\tightitems
      \item Window-specific earnings response coefficients using standardized unexpected earnings.
      \item Covariance between current returns and future earnings innovations.
      \item Horizon-specific price informativeness in the spirit of Bai, Philippon, and Savov.
      \item Full-cycle price accuracy relative to subsequently realized fundamentals.
    \end{itemize}
    \end{block}
  \end{column}
  \begin{column}{0.48\textwidth}
    \begin{block}{Diagnostic measures only}
    \small
    \begin{itemize}\tightitems
      \item Absolute returns and return variance.
      \item Hou-Moskowitz price delay.
      \item Variance ratios and return synchronicity.
      \item Post-earnings-announcement drift.
    \end{itemize}
    \end{block}
  \end{column}
\end{columns}

\begin{exampleblock}{Measurement hierarchy}
Lead with return-fundamental mapping. Use variance, delay, and drift to interpret the timing and persistence of learning, not as substitutes for price informativeness.
\end{exampleblock}

\source{Methodological anchors: Hou and Moskowitz (2005); Bai, Philippon, and Savov (2016); Davila and Parlatore (2026).}
\end{frame}

% -----------------------------------------------------------------------------
\begin{frame}[t]{A window-by-window design separates reallocation from destruction}
\footnotesize
For $w\in\{\text{pre-filing, filing, post-filing, full cycle}\}$, estimate:

\begin{block}{Core specification}
\centering
$R^{w}_{iq}=\alpha_w+\beta_w News_{iq}
+\theta_w\bigl(Restricted_{iq}\times News_{iq}\bigr)
+\Gamma_w X_{iq}+\mu_i+\tau_q+\varepsilon^{w}_{iq}$
\end{block}

\renewcommand{\arraystretch}{1.16}
\begin{tabularx}{\textwidth}{>{\bfseries\RaggedRight\arraybackslash}p{2.50cm} >{\RaggedRight\arraybackslash}X >{\RaggedRight\arraybackslash}X >{\RaggedRight\arraybackslash}X}
\toprule
Scenario & Pre-filing & Filing / post-filing & Full-cycle implication \\
\midrule
Temporal reallocation & $\theta_{pre}<0$ & $\theta_{filing}>0$, little drift & $\sum_w\theta_w\approx0$ \\
Information destruction & $\theta_{pre}<0$ & Filing response does not compensate; later learning persists & $\sum_w\theta_w<0$ \\
Depth / composition only & Little change in return-news coefficients & Amihud or price impact rises & Total fundamental mapping is unchanged \\
\bottomrule
\end{tabularx}

\begin{alertblock}{Interpretive payoff}
\scriptsize If filing-day response rises for the same news, part of the Amihud result is information concentration in the numerator, not simply worse liquidity.
\end{alertblock}
\end{frame}

% -----------------------------------------------------------------------------
\begin{frame}[t]{Essay 1 revision should proceed through three gates}
\begin{columns}[T,onlytextwidth]
  \begin{column}{0.31\textwidth}
    \begin{block}{Gate 1: discovery first stage}
    \small
    Test whether restriction changes indexed-page counts, search ranking, cache recency, discovery latency, or verified crawler consumption.

    \vspace{0.10cm}
    EDGAR downloads are an indirect spillover, not the clean first stage.
    \end{block}
  \end{column}
  \begin{column}{0.31\textwidth}
    \begin{block}{Gate 2: Amihud decomposition}
    \small
    Show event-window and benchmark components separately; add price-impact/depth proxies, breadth, volume, and spread evidence.

    \vspace{0.10cm}
    Avoid claiming a level liquidity collapse that the current data do not show.
    \end{block}
  \end{column}
  \begin{column}{0.31\textwidth}
    \begin{block}{Gate 3: price learning}
    \small
    Link returns to fundamental news across the pre-filing, filing, and post-filing windows.

    \vspace{0.10cm}
    Determine whether information is delayed, destroyed, or simply incorporated with greater price impact.
    \end{block}
  \end{column}
\end{columns}

\begin{exampleblock}{Revised Essay 1 contribution}
\footnotesize The paper would establish a corporate choice over discoverability, document its market consequence, and separate depth, composition, and the relocation of price learning toward mandatory disclosure dates.
\end{exampleblock}
\end{frame}

% -----------------------------------------------------------------------------
\begin{frame}[t]{Essay 2: machine discovery may function as external monitoring}
\begin{block}{Research question}
\large
When a firm restricts machine access to continuously updated public information, does outside monitoring weaken enough to change managerial behavior?
\end{block}

\small
\begin{columns}[T,onlytextwidth]
  \begin{column}{0.48\textwidth}
    \begin{block}{Economic mechanism}
      Automated discovery is an input into scalable anomaly detection. Restrictions raise monitoring costs, reduce the probability that irregular behavior is noticed, and may increase managers' private benefit from opportunism.
    \end{block}
  \end{column}
  \begin{column}{0.48\textwidth}
    \begin{block}{Why this is not Essay 1}
      The outcome is managerial conduct rather than market liquidity. The mechanism operates through detection and discipline, not through the composition of traders at a disclosure event.
    \end{block}
  \end{column}
\end{columns}

\source{Conceptual precedents: Zhu (2019) on data as governance; Heese, Perez-Cavazos, and Peter (2022) on the loss of local monitoring; Diamond and Verrecchia (1991).}
\end{frame}

% -----------------------------------------------------------------------------
\begin{frame}[t,shrink=3]{The monitoring channel must focus on information outside mandatory filings}
\begin{columns}[T,onlytextwidth]
  \begin{column}{0.47\textwidth}
    \begin{exampleblock}{Information plausibly affected}
    \small
    \begin{itemize}\tightitems
      \item Investor-relations page updates.
      \item Product and operating announcements.
      \item Hiring, location, and organizational changes.
      \item Corrections, disappearing pages, and other website changes detectable only through repeated crawling.
    \end{itemize}
    \end{exampleblock}
  \end{column}
  \begin{column}{0.47\textwidth}
    \begin{alertblock}{Information that remains available}
    \small
    \begin{itemize}\tightitems
      \item 10-K, 10-Q, 8-K, and Form 4 filings on EDGAR.
      \item Historical data already collected by institutions.
      \item Bloomberg, direct feeds, and other proprietary channels.
      \item Human access to pages that robots are prohibited from crawling.
    \end{itemize}
    \end{alertblock}
  \end{column}
\end{columns}

\begin{block}{Required exposure measure}
The paper needs a pre-treatment measure of \textbf{website-native monitoring dependence}: how much economically relevant information is updated on the corporate site and not promptly mirrored in EDGAR or major feeds.
\end{block}
\end{frame}

% -----------------------------------------------------------------------------
\begin{frame}[t]{Opportunistic insider trading should be the primary outcome family}
\scriptsize
\renewcommand{\arraystretch}{1.18}
\setlength{\tabcolsep}{0.10cm}
\begin{tabularx}{\textwidth}{>{\bfseries\RaggedRight\arraybackslash}p{2.55cm} >{\RaggedRight\arraybackslash}X >{\RaggedRight\arraybackslash}X}
\toprule
Outcome family & Role in the design & Expected timing \\
\midrule
Opportunistic Form 4 trades & \textbf{Primary}: frequency, trade direction before news, and post-trade abnormal profits & Can move quickly after monitoring weakens \\
Routine Form 4 trades & \textbf{Within-dataset placebo}: mechanically scheduled trades should respond much less & No systematic change \\
Accrual and real earnings management & External validity: behavior that requires managerial discretion and accumulates over time & Slower, quarterly or annual response \\
Reporting discipline & Mechanism extension: NT filings, filing delays, amendments, and restatements & Delayed and potentially asymmetric \\
\bottomrule
\end{tabularx}

\begin{alertblock}{Avoid outcome mining}
\scriptsize Prespecify opportunistic insider trading as the main family. Use accounting choices and reporting discipline only to test generalization and predicted timing.
\end{alertblock}

\source{Classification template: Cohen, Malloy, and Pomorski (2012).}
\end{frame}

% -----------------------------------------------------------------------------
\begin{frame}[t]{Essay 2 faces a sharper endogeneity problem than Essay 1}
\begin{columns}[T,onlytextwidth]
  \begin{column}{0.49\textwidth}
    \begin{alertblock}{Threat}
    \footnotesize
    A firm may restrict crawlers because it anticipates bad news, litigation, website attacks, or opportunistic behavior. Pre-trend tests alone cannot eliminate event-time anticipation.
    \end{alertblock}
    \begin{block}{Design response}
    \footnotesize
    \begin{itemize}\tightitems
      \item Separate active firm-specific changes from vendor/template waves.
      \item Control for website redesigns, security incidents, and disclosure changes.
      \item Compare routine and opportunistic trades within the same filing system.
    \end{itemize}
    \end{block}
  \end{column}
  \begin{column}{0.47\textwidth}
    \begin{block}{Identification target}
    \small
    Reuse the BJS staggered framework only after establishing:
    \begin{enumerate}\tightitems
      \item an actual decline in machine discovery;
      \item no pre-trend in insider trading or reporting behavior;
      \item stronger effects where website-native monitoring should matter;
      \item comparable findings in plausibly passive adoption waves.
    \end{enumerate}
    \end{block}
    \begin{exampleblock}{Interpretation discipline}
    \scriptsize High institutional ownership is heterogeneity, not a perfect placebo: proprietary and web-based monitoring may be complements.
    \end{exampleblock}
  \end{column}
\end{columns}
\end{frame}

% -----------------------------------------------------------------------------
\begin{frame}[t]{Essay 2 opens only if three kill-gates pass}
\begin{columns}[T,onlytextwidth]
  \begin{column}{0.31\textwidth}
    \begin{block}{1. Consumption gate}
      \small Restrictions must reduce an independently measured form of machine discovery or repeated website monitoring.
    \end{block}
  \end{column}
  \begin{column}{0.31\textwidth}
    \begin{block}{2. Pre-trend gate}
      \small Opportunistic trading, bad-news indicators, and reporting behavior must not begin changing before restriction.
    \end{block}
  \end{column}
  \begin{column}{0.31\textwidth}
    \begin{block}{3. Passive-wave gate}
      \small Results must survive separating vendor/template-driven changes from firm-specific active changes.
    \end{block}
  \end{column}
\end{columns}

\vspace{0.20cm}
\begin{exampleblock}{Why the idea remains attractive}
Form 4 and EDGAR reporting outcomes cover small and low-coverage firms that disappear from IBES, 13F, or TAQ-based designs. This aligns the outcome data with the firms for which Essay 1 finds the largest effects.
\end{exampleblock}

\begin{alertblock}{Main risk}
Insider-profit outcomes are noisy and the treatment's causal first stage is currently incomplete. The Form 4 prototype and first-stage pilot should precede full data construction.
\end{alertblock}
\end{frame}

% -----------------------------------------------------------------------------
\begin{frame}[t]{Essay 3 asks whether searchability replaces a lost analyst}
\begin{block}{Research question}
\large
When a broker closure or research-department shutdown removes analyst coverage, does strong pre-existing web discoverability buffer the capital-market damage?
\end{block}

\footnotesize
\begin{columns}[T,onlytextwidth]
  \begin{column}{0.48\textwidth}
    \begin{block}{Why this design is distinct}
      The treatment is an external shock to the human intermediary. Discoverability is measured before the event and enters as a moderator. The paper does not estimate another consequence of the firm's robots decision.
    \end{block}
  \end{column}
  \begin{column}{0.48\textwidth}
    \begin{block}{Why it solves the coverage problem}
      Treated firms have analyst coverage by construction. The outcome panel therefore exists exactly where the coverage-loss treatment occurs, avoiding the empty-cell problem in a robots-to-analyst-production design.
    \end{block}
  \end{column}
\end{columns}

\source{Shock-design precedents: Kelly and Ljungqvist (2012); Derrien and Kecskes (2013). Infrastructure precedents: Chang, Ljungqvist, and Tseng (2023); Guo and Mota (2021).}
\end{frame}

% -----------------------------------------------------------------------------
\begin{frame}[t]{Substitution and complementarity make opposite predictions}
\small
Let $Q$ be a market-quality outcome for which higher values are better. Estimate the effect of analyst loss as a function of pre-event discoverability $D_i$:

\begin{block}{Core interaction}
\centering
$Q_{it}=\alpha_i+\tau_t+\beta Loss_{it}
+\gamma\bigl(Loss_{it}\times D_i^{pre}\bigr)+\Gamma X_{it}+\varepsilon_{it}$
\end{block}

\scriptsize
\renewcommand{\arraystretch}{1.12}
\begin{tabularx}{\textwidth}{>{\bfseries\RaggedRight\arraybackslash}p{2.80cm} >{\RaggedRight\arraybackslash}X >{\RaggedRight\arraybackslash}X}
\toprule
Mechanism & Economic logic & Prediction for $\gamma$ \\
\midrule
Substitution & After the processed analyst channel disappears, investors fall back on searchable raw information. High discoverability reduces the damage. & $\gamma>0$ for market quality; $\gamma<0$ if the outcome is illiquidity \\
Complementarity & Analysts use discoverable raw information as an input. Firms with high discoverability had more productive analyst coverage and lose more when it disappears. & $\gamma<0$ for market quality; $\gamma>0$ if the outcome is illiquidity \\
\bottomrule
\end{tabularx}

\source{Raw vs. processed information: Dugast and Foucault (2018).}
\end{frame}

% -----------------------------------------------------------------------------
\begin{frame}[t]{Predetermined discoverability is useful, but not automatically exogenous}
\begin{columns}[T,onlytextwidth]
  \begin{column}{0.48\textwidth}
    \begin{block}{Preferred moderator}
    \footnotesize
    A pre-event discoverability index combining:
    \begin{itemize}\tightitems
      \item crawler permissions;
      \item indexed-page breadth;
      \item ranking or official-site visibility;
      \item cache/update latency;
      \item share of material information available on the open web.
    \end{itemize}
    \end{block}
  \end{column}
  \begin{column}{0.48\textwidth}
    \begin{alertblock}{What ``predetermined'' does not solve}
    \footnotesize
    Discoverability may correlate with firm size, investor-relations quality, disclosure policy, and the strength of the remaining analyst ecosystem. Measuring it before closure prevents reverse timing but does not make it random.
    \end{alertblock}
  \end{column}
\end{columns}

\begin{exampleblock}{Identification discipline}
\scriptsize Compare firms within the same broker-closure event; include broker-event-by-time structure, baseline coverage controls, balance tests, and moderator bins. Call the result treatment-effect heterogeneity unless discoverability is itself plausibly exogenous.
\end{exampleblock}
\end{frame}

% -----------------------------------------------------------------------------
\begin{frame}[t]{Essay 3 should prioritize information replacement, not Amihud again}
\scriptsize
\renewcommand{\arraystretch}{1.18}
\setlength{\tabcolsep}{0.10cm}
\begin{tabularx}{\textwidth}{>{\bfseries\RaggedRight\arraybackslash}p{2.55cm} >{\RaggedRight\arraybackslash}X >{\RaggedRight\arraybackslash}X}
\toprule
Outcome layer & Measures & Purpose \\
\midrule
Price informativeness & Future-earnings response, price delay, PEAD, horizon-specific informativeness & Primary evidence that open-web information replaces or complements analysts \\
Intermediary response & Entry by other analysts, changes in remaining forecast accuracy or dispersion & Shows whether the analyst ecosystem repairs the loss \\
Information demand & EDGAR or website consumption where valid; media attention & Shows whether investors substitute toward raw sources \\
Firm response & 8-K Item 2.02, earnings releases, and guidance if available & Tests whether firms compensate by supplying more disclosure \\
Investor base and liquidity & Holding breadth, Amihud, and spread proxies & Secondary capital-market consequences \\
\bottomrule
\end{tabularx}

\begin{exampleblock}{Boundary from Essay 1}
\scriptsize Essay 1 studies a firm restriction. Essay 3 asks whether pre-existing discovery absorbs an exogenous loss of human intermediation. Amihud can be secondary, but should not define the contribution.
\end{exampleblock}

\source{Firm disclosure response after coverage loss: Balakrishnan, Billings, Kelly, and Ljungqvist (2014).}
\end{frame}

% -----------------------------------------------------------------------------
\begin{frame}[t]{The broker-shutdown design has one immediate kill-gate: event count}
\begin{columns}[T,onlytextwidth]
  \begin{column}{0.48\textwidth}
    \begin{block}{Two-week feasibility audit}
    \small
    \begin{enumerate}\tightitems
      \item Build a 2016-2024 list of broker closures, mergers, and research-department shutdowns.
      \item Verify events using public news and map affected analysts in IBES Detail.
      \item Count firms, coverage lost, and the low-discoverability interaction cell.
      \item Check whether post-event outcomes remain observable for at least four quarters.
    \end{enumerate}
    \end{block}
  \end{column}
  \begin{column}{0.48\textwidth}
    \begin{alertblock}{Kill criterion}
      \small A large total event count is insufficient. The design needs enough firms that both lose meaningful coverage and differ substantially in pre-event discoverability.
    \end{alertblock}
    \begin{exampleblock}{Fallback sequence}
      \small Expand carefully to broker-level research layoffs or analyst-team departures. If the clean shock remains too thin, transplant the same moderator design to local-newspaper closures rather than forcing a low-powered broker sample.
    \end{exampleblock}
  \end{column}
\end{columns}
\end{frame}

% -----------------------------------------------------------------------------
\begin{frame}[t]{The three-paper architecture is now cleaner than the earlier version}
\scriptsize
\renewcommand{\arraystretch}{1.18}
\setlength{\tabcolsep}{0.10cm}
\begin{tabularx}{\textwidth}{>{\bfseries\RaggedRight\arraybackslash}p{2.15cm} >{\RaggedRight\arraybackslash}X >{\RaggedRight\arraybackslash}X >{\RaggedRight\arraybackslash}X}
\toprule
 & \textbf{Essay 1: Trading} & \textbf{Essay 2: Monitoring} & \textbf{Essay 3: Intermediation} \\
\midrule
Central question & What happens when firms restrict machine discovery? & Does reduced discovery weaken external discipline? & Can discovery capacity replace a lost analyst? \\
Treatment & Robots restriction & Robots restriction & Broker-level coverage loss \\
Primary outcome & Filing-window liquidity and fundamental response timing & Opportunistic insider trading & Price informativeness after analyst loss \\
Distinct mechanism & Depth, composition, and timing of learning & Detection probability and managerial private benefits & Raw vs. processed information substitution \\
Main empirical gate & Actual discoverability first stage & Consumption first stage plus no behavioral pre-trend & Sufficient broker-event-by-discoverability cells \\
Literature community & Microstructure & Governance & Information intermediaries \\
\bottomrule
\end{tabularx}

\begin{block}{Thesis-level contribution}
\scriptsize Corporate information policy is not exhausted by legal publicity. Its effect also depends on machine discoverability and investors' processing capacity.
\end{block}
\end{frame}

% -----------------------------------------------------------------------------
\begin{frame}[t]{Recommended sequence minimizes wasted data work}
\begin{enumerate}
  \item \key{Run the discoverability pilot first.} This is the life-or-death gate for Essay 1, the monitoring channel in Essay 2, and the moderator quality in Essay 3.
  \item \key{Add the price-learning module to Essay 1.} Test temporal reallocation, information destruction, and depth/composition before adding broad outcome families.
  \item \key{Run the broker-event count audit.} A two-week count determines whether Essay 3 can be identified with adequate power.
  \item \key{Prototype Form 4 outcomes.} Measure the incidence and noise of opportunistic trades before committing to a full historical parse.
\end{enumerate}

\vspace{0.18cm}
\begin{alertblock}{Why this order}
The first two pilots answer whether the core friction is measurable and whether the independent-treatment essay is feasible. The Form 4 prototype then determines whether the governance paper has a credible primary outcome.
\end{alertblock}
\end{frame}

% -----------------------------------------------------------------------------
\begin{frame}[t]{Decisions for the supervisor meeting}
\begin{block}{My recommendation}
\large
Keep ``When and Through Which Channel Do Prices Learn?'' inside Essay 1. Develop the governance and intermediary-substitution ideas as the two candidate later essays.
\end{block}

\vspace{0.12cm}
\begin{enumerate}
  \item Is the proposed thesis-level primitive - effective public information depends on discoverability and processing capacity - strong enough to organize all three essays?
  \item For Essay 2, is opportunistic insider trading the right primary governance outcome, conditional on passing the first-stage and pre-trend gates?
  \item For Essay 3, should the immediate priority be a broker-closure event-count audit before building additional outcomes?
  \item What evidence would the committee require before allowing the phrase ``firms actively make themselves harder to find''?
\end{enumerate}

\vfill
\centering
\large\textcolor{MadridBlue}{\textbf{Coherence comes from one economic friction, not from reusing one dataset.}}
\end{frame}

\end{document}
